\documentclass[11pt]{article}
\usepackage[a4paper,margin=25mm]{geometry}
\usepackage{amsmath,amssymb,amsthm}
\usepackage[hidelinks]{hyperref}
\emergencystretch=2em
\newtheorem{theorem}{Theorem}
\newtheorem{lemma}{Lemma}
\newcommand{\R}{\mathbb R}
\newcommand{\Z}{\mathbb Z}
\newcommand{\C}{\mathbb C}
\title{A compact semialgebraic function without exponential Fourier decay\\
\large Disproof of Conjecture 00000008046}
\author{gaochengzhecpu}\date{}
\begin{document}\maketitle
\begin{abstract}
The characteristic function of $[0,1]$ is definable, compactly supported,
and subanalytic. Nevertheless, its Fourier transform has modulus
$1/(\pi(n+1/2))$ at frequency $n+1/2$. This excludes every eventual
bound $Ce^{-c|\xi|}$ with $c>0$, disproving the claimed implication
from subanalyticity to exponential Fourier decay.
\end{abstract}

\section*{The function satisfies the stated hypotheses}
The source concerns definable compactly supported functions and claims
that exponential Fourier decay holds if and only if the function is
subanalytic. No continuity or differentiability hypothesis is imposed.
Consider
\[
 f(x)=\mathbf1_{[0,1]}(x).
\]
Its support is exactly the compact interval $[0,1]$. Its real graph is
\[
 \{(x,y):0\le x\le1,\ y=1\}
 \;\cup\;
 \{(x,y):(x<0\ \text{or}\ x>1),\ y=0\}.
\]
This finite Boolean combination of polynomial conditions makes the
graph semialgebraic, hence definable in the ordered real field and in
the expansions appearing in the source. Polynomial functions are
real analytic, so the displayed description is also semianalytic.
Locally a semianalytic set is subanalytic by the definition using
projections of relatively compact semianalytic sets (restrict to a
bounded coordinate neighborhood and use the identity projection).
Thus the graph, and hence the function in the standard graph sense,
is subanalytic. The discontinuities at $0,1$ do not violate any
hypothesis in the source. Also $f\in L^1(\R)$ since $\int|f|=1$.

\section*{The genuine Fourier integral}
Use the convention
\[
 \widehat f(\xi)=\int_{\R}f(x)e^{-2\pi i x\xi}\,dx.
\]
For $\xi\ne0$, direct integration on the support gives
\[
 \widehat f(\xi)=\int_0^1e^{-2\pi i x\xi}\,dx
       =\frac{e^{-2\pi i\xi}-1}{-2\pi i\xi}.
\]
At $\xi_n=n+1/2$ for $n=0,1,\ldots$, the exponential is $-1$.
Consequently
\[
 \widehat f(\xi_n)=\frac{-i}{\pi\xi_n},\qquad
 |\widehat f(\xi_n)|=\frac1{\pi\xi_n}.
\]
These are exact integral values, not an approximation or an
assumption about a generic Fourier decay model.

\section*{Failure of every exponential bound}
Suppose constants $C>0$, $c>0$, and $R\in\R$ gave
$|\widehat f(\xi)|\le Ce^{-c|\xi|}$ for every $|\xi|\ge R$.
For sufficiently large $n$, the positive frequency $\xi_n$ exceeds
$R$, so the exact values above imply
\[
 1\le\pi C\xi_n e^{-c\xi_n}.
\]
But $x e^{-cx}\to0$ as $x\to+\infty$ for every $c>0$.
The right side therefore tends to zero, a contradiction.
Indeed the same contradiction excludes a bound on just the positive
tail $\xi\ge R$. The alternative Fourier convention
$\int f(x)e^{-i x\xi}\,dx$ only rescales frequency and gives the
same failure. Thus this admissible subanalytic function disproves
the stated ``if'' direction, which suffices to refute the equivalence.

\section*{Correspondence with Lean}
The function \texttt{f} in \texttt{lean/Main.lean} is the actual
interval indicator, viewed in $\C$. The theorem
\texttt{graph\_\allowbreak{}description} proves the displayed finite polynomial
description, including that its imaginary coordinate is zero.
This is the explicit certificate for the semialgebraic and subanalytic
hypotheses. Those general geometric categories are not renamed or
replaced by an unrelated predicate in the Lean file.

The theorem \texttt{compact\_\allowbreak{}support} proves genuine Mathlib compact
support. The theorem \texttt{integral\_\allowbreak{}formula} evaluates
\texttt{Real.fourierIntegral} using its actual Bochner integral,
restriction to $[0,1]$, the equality of closed/open-endpoint integrals
under Lebesgue measure, and the exponential antiderivative.
The theorems \texttt{fourier\_\allowbreak{}at\_\allowbreak{}frequency} and
\texttt{norm\_\allowbreak{}fourier\_\allowbreak{}at\_\allowbreak{}frequency} give the exact half-integer
values. Finally \texttt{no\_\allowbreak{}exponential\_\allowbreak{}bound} quantifies over
arbitrary real $C,R$ and every $c>0$, using the actual exponential
limit and arbitrarily large half-integer frequencies. It is stronger
than the contradiction needed for positive $C$.

\section*{Reproduction}
The project pins Lean 4.19.0 and Mathlib commit
\\\texttt{c44e0c8ee63ca166450922a373c7409c5d26b00b}.\par
From \texttt{lean/}, run \texttt{lake build} and
\texttt{lake env lean -DwarningAsError=true Main.lean}.
On a new machine, \texttt{lake exe cache get} retrieves the official
pinned dependency artifacts. The submission itself is checked in a fresh
build directory. Printed dependencies use only the standard
\texttt{propext}, \texttt{Classical.choice}, and \texttt{Quot.sound}.
No additional axioms, proof placeholders, or numerical oracles are used.
Run \texttt{tectonic main.tex} to reproduce the PDF; no auxiliary numerical
script is needed. Build logs, source hashes, and visual-review records
are included under \texttt{verification/}.
\section*{Conjecture source}
\href{https://github.com/The-Last-Math-Competition/The-Last-Math-Competition/blob/main/conjectures/00000008046.md}{The Last Math Competition, Conjecture 00000008046}. The exact bilingual source and its commit and SHA-256 are preserved in this submission.
\end{document}
